\section{Limitations}
\label{sec:limits}

\paragraph{Exact tables are needed.}
The split requires the answer to every input pair. We obtain it from programs, scene graphs and annotations, which
restricts the analysis to datasets such as CLEVR, GQA, CLOVE and POPE, and even there only to questions well posed on
every image. On GQA, CLOVE and POPE these are existence questions, whose exceptions mostly record which objects
co-occur in which images: basic grounding rather than compositional reasoning. Exact computation is NP-hard in general;
merging identical columns makes it practical for tables with about ten rows, which is a real constraint on table shape.

\paragraph{What the measure is.}
The exception share is a property of a sampled table, not of a task in the abstract. It grows slowly with the number of
images (on CLEVR from 0.079 at 25 images to 0.099 at 1{,}000) and is lower with fewer questions, while rankings of tables
are stable across sizes (Spearman $\ge0.97$; Sec.~\ref{sup:stability}). Comparisons should therefore be made at a fixed
table shape. The share is unique, but individual exception cells are not; 3--6\% of cells are ambiguous across optimal
rules, and cell-level analyses depend on this choice unless restricted to stable labels.

\paragraph{Strength of the VLM result.}
The exception share sets the interaction dependence an accurate model needs, and across four VLMs the EMAP drop follows
it. Most of this relation is what accuracy alone requires: an accuracy-matched null scorer reproduces it, and only
InternVL3.5 relies on interactions clearly beyond the requirement. The comparison with synergy holds within benchmarks
at fixed answer balance, explains little variance, and synergy keeps a signal for one of four models. We tested
zero-shot 7--8B models; larger or fine-tuned models may differ. Our critique of PID concerns deterministic answers on
index-level tables; PID estimators on noisy, high-dimensional features face additional, separate
problems~\cite{liang2023quantifying}.

\paragraph{Scope of the continual-learning results.}
Most continual-learning runs use small fusion networks on frozen features, far from the VLMs that motivate the paper.
The robust findings (forgetting and learning gaps that survive answer matching, difficulty matching and label
ambiguity) are on CLEVR; on GQA and CLOVE answer matching removes most of the effect, and the constructed tasks' second
table has no stable exceptions. For VLMs we tested one model with LoRA, and its pre-registered learning check failed.

\section{Conclusion}
\label{sec:conclusion}
Synergy is widely used to measure how much a multimodal task requires fusion. For deterministic answers it depends
only on per-modality statistics, while the need for non-additive fusion depends on how answers are arranged across
input pairs. Using the classical theory of additive structure, every answer table splits into an additive rule and a
minimum set of exceptions, whose share is the lowest error any additive model can reach and sets a floor on how much an
accurate model must rely on interactions. Four VLMs rely on interactions roughly as much as this floor and their
accuracy require, and one model clearly more. In continual learning, exceptions on compositional tasks are forgotten
faster and learned worse than equally difficult rule cells with the same answer, while standard backward transfer
misses both.

Two recommendations follow. Claims that a benchmark requires fusion should be supported by its exception share, not
by synergy. Continual learners should be evaluated on exceptions separately, against rule cells with the same answer
and difficulty. Natural next steps are estimators of the exception share for data without answer tables, and methods
that learn and keep exceptions without giving up the rule.
